[[1.1]] In the LEP storage ring at CERN, head-on collisions between (equally accelerated) electrons and positrons were produced, such that the total energy in the center of mass was equal to that of the Z boson ( $m_z=91 \mathrm{GeV}$ ). What is the velocity of each particle before the collision? If an electron is accelerated toward a positron at rest, what velocity does it need in order to reach the same center-of-mass total energy?
[[2.1]] The Klein-Gordon equation for a free particle reads

$$
-\frac{\partial^2}{\partial t^2} \phi=-\nabla^2 \phi+m^2 \phi .
$$


Derive the corresponding continuity equation using a similar procedure as in Exercise 1.3.
Verify that the continuity equation can be written as $\partial_\mu j^\mu=0$ where

$$
j_\mu=i\left(\phi^*\left(\partial_\mu \phi\right)-\left(\partial_\mu \phi^*\right) \phi\right) .
$$


What happens to $j_0=\rho$ for a negative energy solution $\phi \sim \exp (i p x)$, and what does this imply for the interpretation of $\rho$ ?
[[3.1]] Given the relativistic invariance of the measure $d^4 k$, show that the integration measure

$$
\frac{d^3 k}{(2 \pi)^3 2 E(\mathbf{k})}
$$

is Lorentz-invariant, provided that $E(\mathbf{k})=\sqrt{\mathbf{k}^2+m^2}$. Use this result to argue that $2 E(\mathbf{k}) \delta^3\left(\mathbf{k}-\mathbf{k}^{\prime}\right)$ is a Lorentz-invariant distribution.
Hint: Start from the Lorentz-invariant expression $\frac{d^4 k}{(2 \pi)^3} \delta\left(k^2-m^2\right) \theta\left(k_0\right)$ and use $\delta\left(x^2-x_0^2\right)=\frac{1}{2|x|}\left[\delta\left(x-x_0\right)+\delta\left(x+x_0\right)\right]$ What is the significance of the $\delta$ and $\theta$ functions above? If you are really keen, you may prove the relation for $\delta\left(x^2-x_0^2\right)$.

[[4.1]]An infinitesimal Lorentz transformation and its inverse can be written as

$$
x^{\prime \alpha}=\left(g^{\alpha \beta}+\epsilon^{\alpha \beta}\right) x_\beta, \quad x^\alpha=\left(g^{\alpha \beta}+\epsilon^{\prime \alpha \beta}\right) x_\beta^{\prime}
$$

where $\epsilon^{\alpha \beta}$ and $\epsilon^{\prime \alpha \beta}$ are infinitesimal. Show from the definition of the inverse that $\epsilon^{\prime \alpha \beta}=- \epsilon^{\alpha \beta}$. Show from the preservation of the norm that $\epsilon^{\alpha \beta}=-\epsilon^{\beta \alpha}$.

[[5.1]]Show that $(\sigma \cdot \mathbf{p})^2=\mathbf{p}^2$. Show that for the solution spinors of the Dirac equation

$$
u_r^{+}(p) u_s(p)=v_r^{+}(p) v_s(p)=2 E \delta_{r s} .
$$


Use the properties of the $\alpha^i, \beta$ (as defined in the Lectures) to show that $\left\{\gamma^\mu, \gamma^\nu\right\}=2 g^{\mu \nu}$.